\documentclass[11pt]{article}
\usepackage[margin=1.05in]{geometry}
\usepackage{amsmath,amssymb,amsthm,mathtools,booktabs,bm}
\usepackage{xcolor}
\usepackage[colorlinks,linkcolor=blue!45!black,citecolor=blue!45!black]{hyperref}

\newtheorem{lemma}{Lemma}[section]
\newtheorem{proposition}[lemma]{Proposition}
\newtheorem{theorem}[lemma]{Theorem}
\newtheorem{corollary}[lemma]{Corollary}
\theoremstyle{definition}
\newtheorem{definition}[lemma]{Definition}
\newtheorem{assumption}[lemma]{Assumption}
\newtheorem{example}[lemma]{Example}
\theoremstyle{remark}
\newtheorem{remark}[lemma]{Remark}
\newtheorem{measurement}[lemma]{Measurement}

\DeclareMathOperator{\Cov}{Cov}
\DeclareMathOperator{\Var}{Var}
\DeclareMathOperator{\spn}{span}
\DeclareMathOperator{\sd}{sd}
\DeclareMathOperator{\Fix}{Fix}
\DeclareMathOperator{\rank}{rank}
\newcommand{\Sph}{\mathbb{S}^{d-1}}
\newcommand{\R}{\mathbb{R}}
\newcommand{\E}{\mathbb{E}}
\newcommand{\T}{\mathsf{T}}
\newcommand{\law}{\mathcal{L}}

\begin{document}

\section{Conventions}\label{sec:conv}

\begin{definition}[second moments]\label{def:conv}
For random vectors $X\in\R^{n}$, $Y\in\R^{m}$,
\[
\Cov(Y,X):=\E\big[(Y-\E Y)(X-\E X)^{\T}\big]\in\R^{m\times n},
\qquad
\Var(X):=\Cov(X,X),
\]
so that $\Cov(Y,X)=\Cov(X,Y)^{\T}$. Unless stated otherwise all random vectors are centred and
$\Cov(Y,X)=\E[YX^{\T}]$.
\end{definition}

\begin{definition}[matrix calculus]\label{def:matcalc}
For a differentiable invertible matrix-valued map $t\mapsto V(t)$,
\begin{equation}\label{eq:dinv}
\frac{d}{dt}V(t)^{-1}=-V(t)^{-1}V'(t)V(t)^{-1},
\end{equation}
obtained by differentiating $VV^{-1}=I$. The identity holds likewise for partial derivatives of a
map $V(p,q)$.
\end{definition}

\begin{definition}[linear ODE, time-ordered exponential]\label{def:texp}
For continuous $K:[0,1]\to\R^{n\times n}$ let $M(t)$ be the unique solution of
\begin{equation}\label{eq:ivp}
\dot M(t)=K(t)M(t),\qquad M(0)=I,
\end{equation}
written $M(t)=\mathcal T\exp\int_{0}^{t}K(s)\,ds$, the limit of the ordered products
$\prod_{i}(I+K(t_i)\Delta t)$ with later times on the left. If $\{K(s)\}_{s}$ is a commuting family
this equals $\exp\int_0^tK$; otherwise it does not.
\end{definition}

\begin{lemma}[Gaussian conditioning]\label{lem:gausscond}
Let $(Y,X)$ be jointly Gaussian and centred with $\Var(X)\succ0$. Then
$\E[Y\mid X=x]=\Cov(Y,X)\Var(X)^{-1}x$.
\end{lemma}

\begin{proof}
Set $L:=\Cov(Y,X)\Var(X)^{-1}$ and $W:=Y-LX$. Then
$\Cov(W,X)=\Cov(Y,X)-L\Var(X)=0$. The pair $(W,X)$ is jointly Gaussian and uncorrelated, hence
independent, so $\E[W\mid X]=\E[W]=0$ and $\E[Y\mid X=x]=Lx$.
\end{proof}

\begin{definition}[symbols]\label{def:symbols}
\begin{center}
\begin{tabular}{ll}
\toprule
$d$ & ambient dimension; $\Sph=\{z\in\R^{d}:\|z\|=1\}$\\
$c$ & condition, $c\sim\mathrm{Unif}(\Sph)$\\
$T_c\Sph$ & $\{v\in\R^d:\langle v,c\rangle=0\}$; $\Pi_z:=I-zz^{\T}$\\
$\xi$ & latent factor, uniform on the unit sphere of $T_c\Sph$\\
$\beta\in[0,1]$ & latent weight; $s:=\sqrt{1-\beta^{2}}$\\
$m$ & $m:=sc+\beta\xi$\\
$\kappa$ & concentration; $A:=A_d(\kappa)=I_{d/2}(\kappa)/I_{d/2-1}(\kappa)$\\
$j$ & leak angle; $z_0:=\cos j\,c+\sin j\,\xi$\\
$z_1$ & $z_1\sim\mathrm{vMF}(m,\kappa)$\\
$\Sigma_0,\Sigma_1,C$ & $\Var(x_0)$, $\Var(x_1)$, $\E[x_0x_1^{\T}]$\\
$S,N$ & $S:=C+C^{\T}$, $N:=\tfrac12(C^{\T}-C)$\\
$P$ & $P:=I-cc^{\T}$\\
\bottomrule
\end{tabular}
\end{center}
\end{definition}

\section{The model}\label{sec:model}

\begin{definition}[generative law]\label{def:model}
Sample in this order:
\[
c\sim\mathrm{Unif}(\Sph);
\qquad
\xi\mid c\sim\mathrm{Unif}\big(\Sph\cap T_c\Sph\big);
\qquad
z_1\mid(c,\xi)\sim\mathrm{vMF}(m,\kappa),\quad m=sc+\beta\xi;
\]
\[
z_0:=\cos j\,c+\sin j\,\xi.
\]
Then $\|m\|^{2}=s^{2}+\beta^{2}=1$ and $\|z_0\|=1$, since $\langle c,\xi\rangle=0$ and
$\|c\|=\|\xi\|=1$.
\end{definition}

\begin{remark}\label{rem:coupling}
$\xi$ is independent of $c$ and enters both $z_0$ (through $\sin j\,\xi$) and $z_1$ (through $m$).
Fixing $\law(z_0\mid c)$ and $\law(z_1\mid c)$ does not determine the joint law of $(z_0,z_1)$
given $c$; the residual freedom is the coupling.
\end{remark}

\begin{definition}[decoupled family]\label{def:rho}
For $\rho\in[0,1]$ let $\xi_s,\xi_b$ be independent and uniform on $\Sph\cap T_c\Sph$ and set
\[
\xi_t:=\frac{\Pi_c\big(\rho\,\xi_s+\sqrt{1-\rho^{2}}\,\xi_b\big)}
{\|\Pi_c(\rho\,\xi_s+\sqrt{1-\rho^{2}}\,\xi_b)\|},
\qquad
m:=s c+\beta\xi_t,
\qquad
z_0:=\cos j\,c+\sin j\,\xi_s,
\qquad
z_1\sim\mathrm{vMF}(m,\kappa).
\]
\end{definition}

\begin{lemma}\label{lem:rhomarg}
For every $\rho\in[0,1]$, $\law(\xi_t\mid c)=\law(\xi_s\mid c)$, hence $\law(z_0\mid c)$ and
$\law(z_1\mid c)$ do not depend on $\rho$. Definition~\ref{def:model} is the case $\rho=1$.
\end{lemma}

\begin{proof}
For $R\in SO(d)$ with $Rc=c$ the pair $(\xi_s,\xi_b)$ is equal in law to $(R\xi_s,R\xi_b)$, and
$\xi_t$ is an $SO(d)$-equivariant function of that pair, so $\law(\xi_t\mid c)$ is invariant under
every such $R$. The only such law on the unit sphere of $T_c\Sph$ is the uniform one. The
statement for $z_1$ follows since $\law(z_1\mid c)$ is determined by $\law(\xi_t\mid c)$.
\end{proof}

\begin{lemma}[first moments]\label{lem:moments}
\begin{equation}\label{eq:m1}
\E[z_1\mid c,\xi]=Am,
\qquad
\E[\xi\mid c]=0,
\qquad
\E[z_1\mid c]=As\,c.
\end{equation}
\end{lemma}

\begin{proof}
The first is the von Mises--Fisher mean. The second holds because $\law(\xi\mid c)$ is invariant
under $\xi\mapsto-\xi$. The third follows by the tower rule,
$\E[z_1\mid c]=\E[A(sc+\beta\xi)\mid c]=Asc+A\beta\E[\xi\mid c]=Asc$.
\end{proof}

\begin{definition}[interpolant, field, objectives]\label{def:flow}
With $\theta:=\arccos\langle z_0,z_1\rangle$,
\[
z_t:=\frac{\sin((1-t)\theta)}{\sin\theta}z_0+\frac{\sin(t\theta)}{\sin\theta}z_1,
\qquad
\dot z_t=\frac{\theta}{\sin\theta}\Big(-\cos((1-t)\theta)z_0+\cos(t\theta)z_1\Big).
\]
The marginal field is
\begin{equation}\label{eq:field}
u(z,t;c):=\E\big[\dot z_t\mid z_t=z,\ c\big].
\end{equation}
Define
\[
\mathcal M_1:\ \min_{\phi}\ \E_{t\sim U[0,1]}\big\|\Pi_{z_t}v_\phi(z_t,t,c)-\dot z_t\big\|^{2},
\qquad
\mathcal M_2:\ \min_{\phi}\ \E_{t\sim U[0,1]}\big\|\Pi_{z_t}v_\phi(z_t,t,c,z_0)-\dot z_t\big\|^{2},
\]
both sampled by $\dot z=\Pi_zv_\phi$ from $z|_{t=0}=z_0$ to $t=1$. The population optimum of
$\mathcal M_1$ is \eqref{eq:field}; that of $\mathcal M_2$ is treated in \S\ref{sec:cond}.
\end{definition}

\section{Reference levels}\label{sec:ref}

\begin{proposition}\label{prop:ref}
Let $z_1'$ be an independent replicate of $z_1$ under the stated conditioning. Then
\begin{align}
\E\langle z_0,z_1\rangle &= A\big(s\cos j+\beta\sin j\big), \label{eq:src}\\
\E\langle m,z_1\rangle &= A, \label{eq:ceil}\\
\E\langle c,z_1\rangle &= As, \label{eq:ceilc}\\
\E\langle z_1',z_1\rangle &= A^{2} &&\hspace{-4em}\text{(replicate given }(c,\xi)), \label{eq:sampx}\\
\E\langle z_1',z_1\rangle &= A^{2}s^{2} &&\hspace{-4em}\text{(replicate given }c). \label{eq:sampc}
\end{align}
\end{proposition}

\begin{proof}
\eqref{eq:src}: $z_0$ is $\sigma(c,\xi)$-measurable, so by \eqref{eq:m1},
\[
\E\langle z_0,z_1\rangle=\E\big\langle z_0,\E[z_1\mid c,\xi]\big\rangle
=A\,\E\langle\cos j\,c+\sin j\,\xi,\ sc+\beta\xi\rangle=A(s\cos j+\beta\sin j).
\]
\eqref{eq:ceil}: $\E\langle m,z_1\rangle=A\langle m,m\rangle=A$.
\eqref{eq:ceilc}: $\E\langle c,z_1\rangle=A\langle c,m\rangle=As$.
\eqref{eq:sampx}: conditionally on $(c,\xi)$ the replicate is independent, so
$\E\langle z_1',z_1\rangle=\|\E[z_1\mid c,\xi]\|^{2}=\|Am\|^{2}=A^{2}$.
\eqref{eq:sampc}: the same step with conditioning on $c$ and $\E[z_1\mid c]=Asc$ gives
$\|Asc\|^{2}=A^{2}s^{2}$.
\end{proof}

\begin{remark}\label{rem:whichopt}
For $\mathcal F$-measurable unit $x$, $\E\langle x,z_1\rangle=\langle x,\E[z_1\mid\mathcal F]\rangle$,
so by Cauchy--Schwarz the maximiser is $x=\E[z_1\mid\mathcal F]/\|\E[z_1\mid\mathcal F]\|$. For
$\mathcal F=\sigma(c,\xi)$ this is $m$, of value $A$; for $\mathcal F=\sigma(c)$ it is $c$, of value
$As$. Hence \eqref{eq:ceil} and \eqref{eq:ceilc} are the maxima of $\E\langle\cdot,z_1\rangle$
under the two information sets, and the functional is maximised by conditional means.
\end{remark}

\begin{corollary}\label{cor:order}
At $d=256$, $A=0.37$, $\beta=0.8$, $j=0.15$,
\[
\underbrace{0.0493}_{\eqref{eq:sampc}}<
\underbrace{0.1369}_{\eqref{eq:sampx}}<
\underbrace{0.2220}_{\eqref{eq:ceilc}}<
\underbrace{0.2637}_{\eqref{eq:src}}<
\underbrace{0.3700}_{\eqref{eq:ceil}}.
\]
\end{corollary}

\begin{measurement}\label{meas:ref}
Monte Carlo with $2\times10^{4}$ draws reproduces \eqref{eq:src}--\eqref{eq:sampc} at
$\beta\in\{0,0.8,0.95\}$ to within $5\times10^{-4}$, against a sampling error of
$4N^{-1/2}=2.8\times10^{-2}$.
\end{measurement}

\section{The flat Gaussian case}\label{sec:flat}

\subsection{Setup}

\begin{assumption}\label{ass:flat}
$(x_0,x_1)$ are jointly Gaussian and centred in $\R^{n}$ with
$\Sigma_0=\Var(x_0)$, $\Sigma_1=\Var(x_1)$, $C=\E[x_0x_1^{\T}]$, and
$J=\begin{psmallmatrix}\Sigma_0&C\\C^{\T}&\Sigma_1\end{psmallmatrix}\succeq0$. The interpolant is
$x_t=(1-t)x_0+tx_1$, so $\dot x_t=x_1-x_0$ is constant in $t$ along each sample path. We assume
$V(t):=\Var(x_t)\succ0$ for all $t\in[0,1]$.
\end{assumption}

\begin{lemma}\label{lem:Vt}
\begin{equation}\label{eq:Vt}
V(t)=(1-t)^{2}\Sigma_0+t(1-t)S+t^{2}\Sigma_1,
\qquad
V(0)=\Sigma_0,\quad V(1)=\Sigma_1.
\end{equation}
\end{lemma}

\begin{proof}
$x_t$ is a linear image of $(x_0,x_1)$, hence Gaussian and centred, with
\[
\E[x_tx_t^{\T}]=(1-t)^{2}\E[x_0x_0^{\T}]+(1-t)t\,\E[x_0x_1^{\T}]+t(1-t)\E[x_1x_0^{\T}]+t^{2}\E[x_1x_1^{\T}].
\]
Substituting $\E[x_1x_0^{\T}]=C^{\T}$ and collecting the middle terms into $t(1-t)(C+C^{\T})$ gives
\eqref{eq:Vt}.
\end{proof}

\subsection{The marginal field}

\begin{lemma}\label{lem:G}
Let $G(t):=\Cov(x_1-x_0,x_t)$. Then
\begin{equation}\label{eq:G}
G(t)=(1-t)\big(C^{\T}-\Sigma_0\big)+t\big(\Sigma_1-C\big).
\end{equation}
\end{lemma}

\begin{proof}
By bilinearity and Definition~\ref{def:conv},
\[
G(t)=\E\big[(x_1-x_0)\big((1-t)x_0+tx_1\big)^{\T}\big]
=(1-t)C^{\T}-(1-t)\Sigma_0+t\Sigma_1-tC. \qedhere
\]
\end{proof}

\begin{lemma}\label{lem:id}
\begin{equation}\label{eq:id}
G(t)=\tfrac12V'(t)+N,\qquad N=\tfrac12\big(C^{\T}-C\big),
\end{equation}
with $N$ skew-symmetric and independent of $t$.
\end{lemma}

\begin{proof}
Differentiating \eqref{eq:Vt},
\[
V'(t)=-2(1-t)\Sigma_0+(1-2t)S+2t\Sigma_1,
\qquad
\tfrac12V'(t)=-(1-t)\Sigma_0+\tfrac{1-2t}{2}\big(C+C^{\T}\big)+t\Sigma_1.
\]
Subtracting from \eqref{eq:G}, the $\Sigma_0$ and $\Sigma_1$ terms cancel, leaving
\[
G(t)-\tfrac12V'(t)
=C^{\T}\Big[(1-t)-\tfrac{1-2t}{2}\Big]+C\Big[-t-\tfrac{1-2t}{2}\Big]
=\tfrac12C^{\T}-\tfrac12C=N,
\]
since $(1-t)-\tfrac{1-2t}{2}=\tfrac12$ and $-t-\tfrac{1-2t}{2}=-\tfrac12$.
\end{proof}

\begin{corollary}\label{cor:idsym}
$G(t)=\tfrac12V'(t)$ for all $t$ if and only if $C=C^{\T}$. For $n=1$ this holds identically.
\end{corollary}

\begin{proposition}\label{prop:linfield}
Under Assumption~\ref{ass:flat},
\begin{equation}\label{eq:K}
u(x,t):=\E[\dot x_t\mid x_t=x]=K(t)x,
\qquad
K(t)=G(t)V(t)^{-1}=\tfrac12V'(t)V(t)^{-1}+NV(t)^{-1}.
\end{equation}
\end{proposition}

\begin{proof}
$(x_1-x_0,x_t)$ is jointly Gaussian and centred, so Lemma~\ref{lem:gausscond} with $Y=x_1-x_0$,
$X=x_t$ gives $u(x,t)=G(t)V(t)^{-1}x$. Substitute \eqref{eq:id}.
\end{proof}

\begin{definition}[endpoint map]\label{def:M}
The probability-flow ODE is $\dot z(t)=K(t)z(t)$, so $z(t)=M(t)z(0)$ with $M$ solving
\eqref{eq:ivp}. Write $M:=M(1)$.
\end{definition}

\begin{remark}\label{rem:fluxnull}
Set $\Lambda(t):=V(t)^{-1}NV(t)^{-1}$. Then $\Lambda^{\T}=-\Lambda$ and
$NV(t)^{-1}=V(t)\Lambda(t)$, so the second term of \eqref{eq:K} is of the form
$V\Lambda$ with $\Lambda$ skew. Proposition~\ref{prop:push} shows this term does not affect the
law of $x_t$; Theorem~\ref{thm:pencil} shows it does affect $M$.
\end{remark}

\subsection{Marginal matching}

\begin{proposition}\label{prop:push}
$M(t)\Sigma_0M(t)^{\T}=V(t)$ for all $t\in[0,1]$; in particular $M\Sigma_0M^{\T}=\Sigma_1$. No
hypothesis on $N$ is required.
\end{proposition}

\begin{proof}
Let $P(t):=M(t)\Sigma_0M(t)^{\T}$, so $P(0)=\Sigma_0=V(0)$, and by \eqref{eq:ivp},
$P'=KP+PK^{\T}$. Evaluate the same expression at $P=V$. Using \eqref{eq:K}, symmetry of $V$ and
$V'$, and $N^{\T}=-N$,
\[
KV=\tfrac12V'+N,
\qquad
K^{\T}=\tfrac12V^{-1}V'-V^{-1}N,
\qquad
VK^{\T}=\tfrac12V'-N,
\]
hence $KV+VK^{\T}=V'$. Therefore $V$ solves $P'=KP+PK^{\T}$ with the same initial condition, and by
uniqueness $P\equiv V$.
\end{proof}

\begin{remark}\label{rem:pushnotenough}
Proposition~\ref{prop:push} does not determine $M$: if $M\Sigma_0M^{\T}=\Sigma_1$ then
$(MQ)\Sigma_0(MQ)^{\T}=\Sigma_1$ for every $Q$ with $Q\Sigma_0Q^{\T}=\Sigma_0$.
\end{remark}

\subsection{The pencil criterion}

\begin{theorem}\label{thm:pencil}
Assume
\[
\textbf{(i)}\ N=0,
\qquad
\textbf{(ii)}\ S\in\spn\{\Sigma_0,\Sigma_1\}.
\]
Then $M$ is a function of $(\Sigma_0,\Sigma_1)$ alone: any two couplings satisfying
\textup{(i)}--\textup{(ii)} with the same marginals produce the same endpoint map.
\end{theorem}

\begin{proof}
\emph{Step 1.} By (ii) write $S=\sigma\Sigma_0+\tau\Sigma_1$. Substituting into \eqref{eq:Vt},
\[
V(t)=p(t)\Sigma_0+q(t)\Sigma_1,
\qquad
p(t)=(1-t)^{2}+\sigma t(1-t),
\qquad
q(t)=t^{2}+\tau t(1-t),
\]
with $(p(0),q(0))=(1,0)$ and $(p(1),q(1))=(0,1)$ for every admissible $(\sigma,\tau)$. The coupling
moves the curve $\gamma(t)=(p(t),q(t))$ but not its endpoints.

\emph{Step 2.} On $\Omega:=\{(p,q)\in\R^{2}:p\Sigma_0+q\Sigma_1\succ0\}$ define
\[
V(p,q):=p\Sigma_0+q\Sigma_1,
\qquad
A(p,q):=\tfrac12\Sigma_0V^{-1},
\qquad
B(p,q):=\tfrac12\Sigma_1V^{-1}.
\]
By (i), $K(t)=\tfrac12V'(t)V(t)^{-1}$ with $V'(t)=p'(t)\Sigma_0+q'(t)\Sigma_1$, hence
$K(t)=p'(t)A(\gamma(t))+q'(t)B(\gamma(t))$. Thus \eqref{eq:ivp} is the pullback along $\gamma$ of
the matrix-valued $1$-form $\omega:=A\,dp+B\,dq$, and $M$ is the parallel transport of $\omega$
along $\gamma$ from $(1,0)$ to $(0,1)$. The set $\Omega$ is convex by linearity of $V$ in $(p,q)$,
hence simply connected.

\emph{Step 3.} If transport is path-independent on $\Omega$ there is $M(p,q)$ with
$\partial_pM=AM$, $\partial_qM=BM$. Cross-differentiating,
\[
\partial_q\partial_pM=(\partial_qA)M+ABM,
\qquad
\partial_p\partial_qM=(\partial_pB)M+BAM,
\]
so for invertible $M$ equality of mixed partials is equivalent to
\begin{equation}\label{eq:frob}
\partial_qA-\partial_pB+[A,B]=0,
\end{equation}
which by the Frobenius theorem is also sufficient on the simply connected $\Omega$.

\emph{Step 4.} From $\partial_pV=\Sigma_0$, $\partial_qV=\Sigma_1$ and \eqref{eq:dinv},
\[
\partial_qA=-\tfrac12\Sigma_0V^{-1}\Sigma_1V^{-1}=-\big(\tfrac12\Sigma_0V^{-1}\big)\big(\Sigma_1V^{-1}\big)=-2AB,
\qquad
\partial_pB=-2BA,
\]
whence $\partial_qA-\partial_pB+[A,B]=-2AB+2BA+AB-BA=-[A,B]$, so \eqref{eq:frob} holds if and only
if $[A,B]=0$.

\emph{Step 5.} By construction $pA+qB=\tfrac12(p\Sigma_0+q\Sigma_1)V^{-1}=\tfrac12I$. For $q\neq0$,
$B=\tfrac{1}{2q}I-\tfrac{p}{q}A$, so
\[
[A,B]=\Big[A,\tfrac{1}{2q}I\Big]-\tfrac{p}{q}[A,A]=0.
\]
The set $\{q=0\}$ is null in $\Omega$ and $[A,B]$ is continuous, so $[A,B]\equiv0$ on $\Omega$.

\emph{Step 6.} By Steps 3--5 transport is path-independent on $\Omega$; by Step 1 the endpoints of
$\gamma$ are $(1,0)$ and $(0,1)$ independently of the coupling, with $V(1,0)=\Sigma_0$ and
$V(0,1)=\Sigma_1$. Hence $M$ depends only on $(\Sigma_0,\Sigma_1)$.
\end{proof}

\begin{remark}\label{rem:notused}
Step 5 does not use $[\Sigma_0,\Sigma_1]=0$. Commutation of the marginals is not required for
coupling-independence; only (i) and (ii) are.
\end{remark}

\begin{corollary}\label{cor:brenier}
If in addition $[\Sigma_0,\Sigma_1]=0$ then $M=\Sigma_1^{1/2}\Sigma_0^{-1/2}$, the Brenier map
between $\mathcal N(0,\Sigma_0)$ and $\mathcal N(0,\Sigma_1)$.
\end{corollary}

\begin{proof}
By Theorem~\ref{thm:pencil} evaluate $M$ along the segment $p(u)=1-u$, $q(u)=u$ in $\Omega$, so
$V(u)=(1-u)\Sigma_0+u\Sigma_1$ and $K=\tfrac12(\Sigma_1-\Sigma_0)V(u)^{-1}$. Since $\Sigma_0$ and
$\Sigma_1$ commute they are simultaneously diagonalisable, so $\{V(u),V'(u)\}_u$ is a commuting
family, the time-ordering of Definition~\ref{def:texp} is vacuous, and
\[
M=\exp\Big(\tfrac12\int_0^1V'V^{-1}\,du\Big)
=\exp\Big(\tfrac12\big[\log\Sigma_1-\log\Sigma_0\big]\Big)=\Sigma_1^{1/2}\Sigma_0^{-1/2}. \qedhere
\]
\end{proof}

\begin{corollary}\label{cor:notbayes}
For $n=1$ with means restored, $M:x_0\mapsto\mu_1+\tfrac{\sigma_1}{\sigma_0}(x_0-\mu_0)$, whereas
$\E[x_1\mid x_0]=\mu_1+\rho\tfrac{\sigma_1}{\sigma_0}(x_0-\mu_0)$ with
$\rho=\Cov(x_0,x_1)/(\sigma_0\sigma_1)$. The two coincide only at $\rho=1$.
\end{corollary}

\subsection{Failure off the pencil}

\begin{proposition}\label{prop:obstruction}
Drop hypothesis \textup{(ii)} and set $V(p,r,q)=p\Sigma_0+rS+q\Sigma_1$ on the corresponding
positivity domain, with $A=\tfrac12\Sigma_0V^{-1}$, $R=\tfrac12SV^{-1}$, $B=\tfrac12\Sigma_1V^{-1}$.
Then $pA+rR+qB=\tfrac12I$ and
\[
[A,B]=-\frac{r}{q}[A,R].
\]
The integrability condition \eqref{eq:frob} in the $(p,q)$ directions therefore holds if and only
if $r[A,R]=0$, and hypothesis \textup{(ii)} is the case $r\equiv0$.
\end{proposition}

\begin{proof}
$pA+rR+qB=\tfrac12VV^{-1}=\tfrac12I$ as in Step 5. Solving for $B$ and commuting with $A$,
\[
q[A,B]=\Big[A,\ \tfrac12I-pA-rR\Big]=-p[A,A]-r[A,R]=-r[A,R].
\]
Steps 3--4 of Theorem~\ref{thm:pencil} apply verbatim in the $(p,q)$ directions, reducing
\eqref{eq:frob} to $[A,B]=0$.
\end{proof}

\begin{remark}\label{rem:obstrscope}
Proposition~\ref{prop:obstruction} establishes failure of a sufficient condition, not
coupling-dependence. The latter is exhibited by Example~\ref{ex:arms}\,(c)--(d).
\end{remark}

\subsection{Numerical verification}

\begin{measurement}\label{meas:id}
Over $300$ random admissible triples $(\Sigma_0,\Sigma_1,C)$ with $n\in\{1,\dots,5\}$ and
$t\sim U(0.05,0.95)$,
\[
\max\big\|G(t)-\tfrac12V'(t)-\tfrac12(C^{\T}-C)\big\|_{\max}=3.6\times10^{-15}.
\]
\end{measurement}

\begin{example}\label{ex:arms}
Let $n=3$, $\Sigma_0=\mathrm{diag}(1,\ 0.4,\ 0.15)$, $D=\mathrm{diag}(0.8,\ 0.5,\ 0.25)$, let
$Q,Q'\in SO(3)$ be fixed and generic, let $E=Q'\mathrm{diag}(1,\ 0.6,\ 0.3)Q'^{\T}$, and let
$W=-W^{\T}$ be a fixed generic skew matrix. For $a\in\{0,0.1,0.2,0.3,0.45,0.6\}$ consider the four
families
\[
\begin{array}{llll}
\text{(a)} & \Sigma_1=D, & C=a\,\mathrm{diag}\big(\sqrt{(\Sigma_0)_{ii}(\Sigma_1)_{ii}}\big), \\[2pt]
\text{(b)} & \Sigma_1=QDQ^{\T}, & C=\tfrac{a}{2}\big(\Sigma_0+\Sigma_1\big), \\[2pt]
\text{(c)} & \Sigma_1=D, & C=0.35\,a\,E, \\[2pt]
\text{(d)} & \Sigma_1=QDQ^{\T}, & C=0.15\big(\Sigma_0+\Sigma_1\big)+a\,W,
\end{array}
\]
restricted to those $a$ for which $J\succeq0$. Write
$\delta:=\max_a\|M(a)-M(0)\|_F/\|M(0)\|_F$.
\end{example}

\begin{measurement}\label{meas:arms}
For Example~\ref{ex:arms}, integrating \eqref{eq:ivp} on a uniform grid of $4000$ steps:
\begin{center}
\begin{tabular}{ccccccc}
\toprule
family & $C=C^{\T}$ & $[\Sigma_0,\Sigma_1]=0$ & $S\in\spn\{\Sigma_0,\Sigma_1\}$ & $\delta$ &
$\|M-\Sigma_1^{1/2}\Sigma_0^{-1/2}\|_{\mathrm{rel}}$ & $\|M\Sigma_0M^{\T}-\Sigma_1\|_{\max}$\\
\midrule
(a) & yes & yes & yes & $6.3\times10^{-5}$ & $1.0\times10^{-4}$ & $1.2\times10^{-4}$\\
(b) & yes & no  & yes & $6.3\times10^{-5}$ & $3.3\times10^{-2}$ & $2.3\times10^{-4}$\\
(c) & yes & yes & no  & $7.6\times10^{-3}$ & $1.1\times10^{-2}$ & $1.0\times10^{-4}$\\
(d) & no  & no  & yes & $9.5\times10^{-1}$ & $3.2\times10^{-1}$ & $1.5\times10^{-4}$\\
\bottomrule
\end{tabular}
\end{center}
The discretisation floor is $10^{-4}$, read from the last column.
\end{measurement}

\begin{remark}\label{rem:armsread}
Family (b) satisfies (i)--(ii) with $[\Sigma_0,\Sigma_1]\neq0$ and is invariant, confirming
Remark~\ref{rem:notused}; it departs from the Brenier map, confirming that
Corollary~\ref{cor:brenier} requires the extra hypothesis. Families (c) and (d) violate (ii) and
(i) respectively and both move, so neither hypothesis alone suffices. In every family the last
column is at the discretisation floor, confirming Proposition~\ref{prop:push} independently of
$N$.
\end{remark}

\subsection{Degenerate sources}

\begin{corollary}\label{cor:degen}
In the scalar case, as $\Var(x_0)\to0$ the endpoint map degenerates to the constant
$x_0\mapsto\mu_1$.
\end{corollary}

\begin{proof}
By Corollary~\ref{cor:notbayes} the gain $\sigma_1/\sigma_0$ multiplies $(x_0-\mu_0)$, which
vanishes identically when $\Var(x_0)=0$.
\end{proof}

\begin{remark}\label{rem:degenmech}
A Lipschitz map cannot increase Hausdorff dimension, hence cannot push a measure onto one of
strictly larger support dimension. When $\law(x_0)$ is degenerate no flow with Lipschitz field is
marginal-preserving, and Assumption~\ref{ass:flat} fails at $t=0$.
\end{remark}

\begin{definition}[regularised family]\label{def:sigfam}
For $\sigma>0$ replace $z_0$ of Definition~\ref{def:model} by
\[
z_0^{(\sigma)}:=\cos\varphi_0\,c+\sin\varphi_0\,\xi,
\qquad
\varphi_0=j+\sigma\varepsilon,
\qquad
\varepsilon\sim\mathcal N(0,1)\ \text{independent of}\ (c,\xi,z_1).
\]
\end{definition}

\begin{lemma}\label{lem:sigrank}
For every $\sigma>0$, $\Var(\langle z_0^{(\sigma)},c\rangle\mid c)>0$ and the support of
$\law(z_0^{(\sigma)}\mid c)$ has dimension $d-1$. For $\sigma=0$ both statements fail: the variance
is $0$ and the support has dimension $d-2$.
\end{lemma}

\begin{proof}
$\langle z_0^{(\sigma)},c\rangle=\cos(j+\sigma\varepsilon)$, which is non-constant for $\sigma>0$
and equals $\cos j$ identically for $\sigma=0$. The support is
$\{\cos\varphi\,c+\sin\varphi\,\xi\}$ over $\varphi$ in an interval and $\xi$ in a $(d-2)$-sphere,
of dimension $(d-2)+1$ for $\sigma>0$ and $d-2$ for $\sigma=0$.
\end{proof}

\begin{proposition}\label{prop:sigmean}
In the scalar Gaussian model with $\Var(x_0)=\sigma^{2}>0$, the endpoint map is
$x_0\mapsto\mu_1+(\sigma_1/\sigma)(x_0-\mu_0)$, so $\E[M(x_0)]=\mu_1$ for every $\sigma>0$ and the
limit $\sigma\downarrow0$ is the constant map to $\mu_1$.
\end{proposition}

\begin{proof}
Corollaries~\ref{cor:notbayes} and \ref{cor:degen}, with $\E[x_0-\mu_0]=0$.
\end{proof}

\begin{measurement}\label{meas:sigflat}
For Definition~\ref{def:sigfam} at $d=256$, $A=0.37$, $\beta=0.8$, $j=0.15$, with second moments
estimated from $1.2\times10^{5}$ draws and the endpoint map computed in the invariant algebra
$\spn\{cc^{\T},P\}$:
\begin{center}
\begin{tabular}{ccccc}
\toprule
$\sigma$ & $\sd(\langle z_0,c\rangle)$ & $\Var(\langle z_0,c\rangle)$ & gain & $\E\langle z_1,c\rangle$\\
\midrule
$0$     & $5.6\times10^{-16}$ & $7.1\times10^{-24}$ & $\infty$ & $0.22241$\\
$0.005$ & $7.5\times10^{-4}$  & $5.6\times10^{-7}$  & $73.65$  & $0.22216$\\
$0.01$  & $1.5\times10^{-3}$  & $2.2\times10^{-6}$  & $36.84$  & $0.22213$\\
$0.03$  & $4.5\times10^{-3}$  & $2.0\times10^{-5}$  & $12.28$  & $0.22212$\\
$0.06$  & $9.3\times10^{-3}$  & $8.6\times10^{-5}$  & $5.983$  & $0.22210$\\
$0.10$  & $1.6\times10^{-2}$  & $2.7\times10^{-4}$  & $3.378$  & $0.22163$\\
$0.20$  & $4.0\times10^{-2}$  & $1.6\times10^{-3}$  & $1.382$  & $0.22204$\\
$0.40$  & $1.2\times10^{-1}$  & $1.4\times10^{-2}$  & $0.4675$ & $0.22185$\\
\bottomrule
\end{tabular}
\end{center}
The gain agrees with $\sd(\langle z_1,c\rangle)/\sd(\langle z_0,c\rangle)$ to a maximum relative
error of $7.1\times10^{-3}$, and $\E\langle z_1,c\rangle$ agrees with $As=0.2220$ to
$4.1\times10^{-4}$ throughout. The pencil determinant $a_0b_1-a_1b_0$ is nonzero at every
$\sigma$, including $\sigma=0$, and changes sign between $\sigma=0.03$ and $\sigma=0.06$.
\end{measurement}

\begin{remark}\label{rem:sigread}
Hypotheses (i) and (ii) of Theorem~\ref{thm:pencil} hold at every $\sigma$ including $\sigma=0$.
Only Assumption~\ref{ass:flat} fails at $\sigma=0$, at the single point $t=0$ and in the single
direction $c$.
\end{remark}

\subsection{Quantitative form}

Theorem~\ref{thm:pencil} is a dichotomy. The following states that it is also a stability
statement: the endpoint map responds to a violation at first order, with a response operator that
annihilates exactly the directions the theorem permits.

\begin{proposition}[first-order response]\label{prop:stab}
Let $C_0$ satisfy \textup{(i)} and \textup{(ii)} of Theorem~\ref{thm:pencil}, let
$E\in\R^{n\times n}$, and put $C_\varepsilon:=C_0+\varepsilon E$. Write $V_0,G_0,K_0$ for the data
of $C_0$ and let $\Phi$ solve $\Phi'=K_0\Phi$, $\Phi(0)=I$, so that $\Phi(t)=M(t)$ for $C_0$ and
$M_0=\Phi(1)$. Then there is $\varepsilon_0>0$ such that
$\varepsilon\mapsto M(C_\varepsilon)$ is real-analytic on $(-\varepsilon_0,\varepsilon_0)$ and
\begin{equation}\label{eq:resp}
M(C_\varepsilon)=M_0+\varepsilon\,\mathcal D(E)+O(\varepsilon^{2}),
\qquad
\mathcal D(E)=M_0\int_0^1\Phi(t)^{-1}K_E(t)\,\Phi(t)\,dt,
\end{equation}
\begin{equation}\label{eq:KE}
K_E(t)=\big[(1-t)E^{\T}-tE\big]V_0(t)^{-1}
-G_0(t)V_0(t)^{-1}\big[t(1-t)(E+E^{\T})\big]V_0(t)^{-1}.
\end{equation}
Moreover $\mathcal D(E)=0$ whenever $E=E^{\T}$ and $E\in\spn\{\Sigma_0,\Sigma_1\}$.
\end{proposition}

\begin{proof}
\emph{Analyticity.} By \eqref{eq:Vt}, $V_\varepsilon(t)=V_0(t)+\varepsilon t(1-t)(E+E^{\T})$ is
affine in $\varepsilon$. Since $V_0(t)\succ0$ on the compact $[0,1]$, its least eigenvalue is
bounded below by some $\lambda>0$, so there is $\varepsilon_0>0$ with $V_\varepsilon(t)\succ0$ for
all $t\in[0,1]$ and $|\varepsilon|<\varepsilon_0$. On that set $V_\varepsilon(t)^{-1}$ is
real-analytic in $\varepsilon$ and continuous in $t$, and $G_\varepsilon(t)$ is affine in
$\varepsilon$ by \eqref{eq:G}; hence $K_\varepsilon(t)=G_\varepsilon(t)V_\varepsilon(t)^{-1}$ is
real-analytic in $\varepsilon$, jointly continuous, and uniformly bounded on
$[0,1]\times(-\varepsilon_0,\varepsilon_0)$. Analytic dependence of the solution of the linear
initial value problem \eqref{eq:ivp} on a parameter gives the first claim.

\emph{The formula.} Differentiating \eqref{eq:ivp} in $\varepsilon$ at $\varepsilon=0$ and writing
$D(t):=\partial_\varepsilon M_\varepsilon(t)|_{\varepsilon=0}$,
\[
D'=K_0D+K_E\Phi,\qquad D(0)=0,
\]
since $M_\varepsilon(t)|_{\varepsilon=0}=\Phi(t)$. Differentiating $G_\varepsilon$ and
$V_\varepsilon$ in $\varepsilon$ and applying $\partial_\varepsilon(V^{-1})=-V^{-1}(\partial_\varepsilon V)V^{-1}$
gives \eqref{eq:KE}. Variation of constants yields
$D(t)=\Phi(t)\int_0^t\Phi(s)^{-1}K_E(s)\Phi(s)\,ds$; evaluate at $t=1$.

\emph{The kernel.} If $E=E^{\T}$ and $E\in\spn\{\Sigma_0,\Sigma_1\}$ then $C_\varepsilon$ is
symmetric, so (i) holds, and $C_\varepsilon+C_\varepsilon^{\T}=2C_0+2\varepsilon E\in
\spn\{\Sigma_0,\Sigma_1\}$, so (ii) holds. By Theorem~\ref{thm:pencil}, $M(C_\varepsilon)=M_0$ for
every such $\varepsilon$, and the derivative vanishes.
\end{proof}

\begin{corollary}\label{cor:stab}
With $\|E\|=1$, $\ \|M(C_\varepsilon)-M_0\|=\varepsilon\|\mathcal D(E)\|+O(\varepsilon^{2})$. The
displacement of the endpoint map is therefore proportional to the size of the violation at first
order, with constant $\|\mathcal D(E)\|$ depending only on the components of $E$ that violate
\textup{(i)} or \textup{(ii)}.
\end{corollary}

\begin{measurement}\label{meas:stab}
Six dimensions $n\in\{2,4,8,16,32,64\}$, six systems each, $2048$ random pairs of marginals per
cell and thirty-six couplings per pair, float64, second-order integration at $2000$ steps. Across
the $73{,}728$ systems satisfying (i)--(ii) the relative displacement has ninetieth percentile
$4.06\times10^{-8}$, against a median of at least $9.78\times10^{-3}$ once a hypothesis is
violated at relative size $0.1$: a separation of $2.4\times10^{5}$. Invariance is unaffected by
$[\Sigma_0,\Sigma_1]$, holding at normalised commutator up to $0.199$, as
Remark~\ref{rem:notused} requires. Regressing $\log_{10}$ displacement on $\log_{10}$ violation,
$180$ points per family, gives slope $0.9828$ with residual $0.0349$ for violations of (ii) and
slope $0.9419$ with residual $0.0563$ for violations of (i): both unity to within the fit, as
Corollary~\ref{cor:stab} requires. At unit relative violation the two responses are $0.647$ and
$0.0982$, a ratio of $6.585$. That ratio is the quotient of the two constants
$\|\mathcal D(E)\|$; it is measured, not predicted.
\end{measurement}

\subsection{The loss floor}

\begin{definition}\label{def:floor}
$\mathcal F(C):=\displaystyle\int_0^1\operatorname{Tr}\Var\big(\dot x_t\mid x_t\big)\,dt$, the
value of the objective $\mathcal M_1$ of Definition~\ref{def:flow} at its population optimum.
\end{definition}

\begin{lemma}\label{lem:floor}
Under Assumption~\ref{ass:flat},
\begin{equation}\label{eq:floor}
\mathcal F(C)=\int_0^1\operatorname{Tr}\Big(\Sigma_0+\Sigma_1-C-C^{\T}
-G(t)V(t)^{-1}G(t)^{\T}\Big)\,dt .
\end{equation}
\end{lemma}

\begin{proof}
$\dot x_t=x_1-x_0$ with $\Var(x_1-x_0)=\Sigma_0+\Sigma_1-C-C^{\T}$. In the proof of
Lemma~\ref{lem:gausscond}, with $Y=x_1-x_0$ and $X=x_t$, the residual $W=Y-LX$ is independent of
$X$, so $\Var(Y\mid X)=\Var(W)=\Var(Y)-\Cov(Y,X)\Var(X)^{-1}\Cov(X,Y)$, which is constant in $x$.
Substitute $\Cov(Y,X)=G(t)$ from Lemma~\ref{lem:G} and $\Var(X)=V(t)$, then integrate.
\end{proof}

\begin{proposition}\label{prop:decouple}
The loss floor is not a function of the endpoint map. Explicitly, take $n=1$ and
$\Sigma_0=\Sigma_1=1$. Every $c\in[-1,1]$ satisfies \textup{(i)} and \textup{(ii)} and gives
$M=1$, while
\begin{equation}\label{eq:F1}
\mathcal F(c)=\int_0^1\Big(2(1-c)-\frac{(1-2t)^{2}(1-c)^{2}}{1-2t(1-t)(1-c)}\Big)dt,
\qquad
\mathcal F(1)=0,\quad \mathcal F(0)=\frac{\pi}{2}.
\end{equation}
Consequently there is no function $\varphi$ with $\varphi(0)=0$ such that
$|\mathcal F(C_1)-\mathcal F(C_2)|\le\varphi\big(\|M(C_1)-M(C_2)\|\big)$ for all admissible
couplings with common marginals.
\end{proposition}

\begin{proof}
For $n=1$ every $C$ is symmetric, so (i) holds; and $\spn\{\Sigma_0,\Sigma_1\}=\R$, so (ii)
holds. Positive semidefiniteness of $J$ is $|c|\le1$. By Corollary~\ref{cor:idsym},
$G=\tfrac12V'$, so $K=\tfrac12V'V^{-1}$ and
$M=\exp\big(\tfrac12\int_0^1V'/V\,dt\big)=\sqrt{V(1)/V(0)}=\sqrt{\Sigma_1/\Sigma_0}=1$,
independently of $c$. By \eqref{eq:Vt}, $V(t)=1-2t(1-t)(1-c)$, hence
$G(t)=\tfrac12V'(t)=-(1-2t)(1-c)$, and $\Var(x_1-x_0)=2-2c$; substituting into \eqref{eq:floor}
gives \eqref{eq:F1}. At $c=1$ both terms vanish. At $c=0$, $V=2t^{2}-2t+1$ and
$(1-2t)^{2}=2V-1$, so the integrand is $2-(2V-1)/V=1/V$, and with $u=t-\tfrac12$,
\[
\int_0^1\frac{dt}{2t^{2}-2t+1}=\int_{-1/2}^{1/2}\frac{du}{2u^{2}+\tfrac12}
=\Big[\arctan(2u)\Big]_{-1/2}^{1/2}=\frac{\pi}{2}.
\]
For the last claim take $c=1$ and $c=0$: the endpoint maps coincide, so the right side is
$\varphi(0)=0$, while the left side is $\pi/2$.
\end{proof}

\begin{measurement}\label{meas:decouple}
Forty-seven random systems in dimensions $3$ and $6$, coupling strength swept along arms
satisfying (i)--(ii) and along arms violating them. On the satisfying arms the endpoint map moves
by a median $3.17\times10^{-4}$ while $\mathcal F$ spans a median factor $2.23$, range
$[2.14,2.41]$. On the violating arms the map moves by a median $6.35\times10^{-1}$ while
$\mathcal F$ spans a median factor $1.23$, range $[1.07,1.99]$.
\end{measurement}

\begin{remark}\label{rem:floorread}
The spread of $\mathcal F$ is therefore larger on the couplings that are provably equivalent at
the endpoint than on those that are not. Ranking couplings by the value of the objective at its
optimum orders them by a quantity that is not merely uninformative about $M$ but ordered against
it. Reports that the coupling changes the attainable loss by orders of magnitude are consistent
with Theorem~\ref{thm:pencil} and do not bear on it: by Proposition~\ref{prop:decouple} the two
quantities are different functions of $C$, and only $M$ is realised at sampling time.
\end{remark}

\section{The sphere}\label{sec:sphere}

\subsection{Symmetry}

\begin{definition}\label{def:stab}
Let $H_c:=\{R\in SO(d):Rc=c\}$, acting on $c^{\perp}\cong\R^{d-1}$ as $SO(d-1)$. For
$z\notin\{\pm c\}$ let $H_{c,z}:=\{R\in H_c:Rz=z\}$, acting on $\{c,z\}^{\perp}\cong\R^{d-2}$ as
$SO(d-2)$.
\end{definition}

\begin{lemma}\label{lem:invlaw}
For every $R\in H_c$, $(Rz_0,Rz_1)\mid c\overset{d}{=}(z_0,z_1)\mid c$.
\end{lemma}

\begin{proof}
$\law(\xi\mid c)$ is uniform on the unit sphere of $c^{\perp}$, which is $H_c$-invariant, so
$R\xi\overset{d}{=}\xi$. Given $\xi$, $Rm=R(sc+\beta\xi)=sc+\beta R\xi$ since $Rc=c$. The vMF
density at $z$ depends on $(z,m)$ only through $\langle z,m\rangle$ and
$\langle Rz,Rm\rangle=\langle z,m\rangle$, so $Rz_1\sim\mathrm{vMF}(Rm,\kappa)$. Finally
$Rz_0=\cos j\,c+\sin j\,R\xi$, so $(Rz_0,Rz_1)$ is generated by Definition~\ref{def:model} from
$(c,R\xi)$.
\end{proof}

\begin{lemma}\label{lem:equiv}
$u(Rz,t;c)=R\,u(z,t;c)$ for all $R\in H_c$.
\end{lemma}

\begin{proof}
The interpolant is built from $\langle z_0,z_1\rangle$ and linear combinations of $z_0,z_1$, both
$SO(d)$-equivariant, so $Rz_t(z_0,z_1)=z_t(Rz_0,Rz_1)$ and likewise for $\dot z_t$. Hence
\[
u(Rz,t;c)=\E[\dot z_t\mid z_t=Rz,c]=\E[R\dot z_t\mid Rz_t=Rz,c]=R\,\E[\dot z_t\mid z_t=z,c],
\]
the middle equality using Lemma~\ref{lem:invlaw}.
\end{proof}

\begin{proposition}\label{prop:2plane}
Let $d\geq4$ and $z\notin\{\pm c\}$. Then $u(z,t;c)\in\spn\{z,c\}$.
\end{proposition}

\begin{proof}
For $R\in H_{c,z}$, $Rz=z$, so Lemma~\ref{lem:equiv} gives $u(z,t;c)=R\,u(z,t;c)$; the vector is
fixed by every element of $H_{c,z}$. Decompose $\R^{d}=\spn\{c,z\}\oplus\{c,z\}^{\perp}$. The
group $H_{c,z}$ acts trivially on the first summand and as $SO(d-2)$ on the second; for
$d-2\geq2$ the only vector of $\R^{d-2}$ fixed by all of $SO(d-2)$ is $0$. Hence
$\Fix(H_{c,z})=\spn\{c,z\}$.
\end{proof}

\begin{corollary}\label{cor:azimuth}
The trajectory of $\dot z=u(z,t;c)$ started at $z_0$ remains in $\Pi:=\spn\{c,z_0\}=\spn\{c,\xi\}$.
Writing $z_t=\cos\varphi_t\,c+\sin\varphi_t\,\xi$ with $\varphi_0=j$, the direction $\xi$ is
constant along the trajectory and only $\varphi_t$ evolves.
\end{corollary}

\begin{proof}
By Proposition~\ref{prop:2plane}, $\dot z\in\spn\{z,c\}$; if $z\in\Pi$ and $c\in\Pi$ then
$\dot z\in\Pi$, so $\Pi$ is invariant and $z_0\in\Pi$. With $\|z_t\|=1$ the trajectory lies on the
great circle $\Pi\cap\Sph$, covered by the stated parameterisation while $\sin\varphi_t\neq0$.
\end{proof}

\begin{measurement}\label{meas:azimuth}
Let $\Sigma:=\E\big[\langle x_\infty,c\rangle^{2}+\langle x_\infty,\xi\rangle^{2}\big]$, which
equals $1$ if and only if $x_\infty\in\Pi$ almost surely. Over $40$ parameter settings
$(\sigma,\rho)$ with $\sigma\in\{0,0.005,0.01,0.03,0.06,0.1,0.2,0.4\}$ and
$\rho\in\{0,0.25,0.5,0.75,1\}$, the population field of Definition~\ref{def:flow} attains
$\Sigma=1.000$ in all $40$ settings. A network trained on $\mathcal M_1$ attains
$\Sigma=0.859$ at $\sigma=0$, decreasing to $0.609$ at $\sigma=0.4$.
\end{measurement}

\subsection{The landing law}

\begin{lemma}\label{lem:polar}
$\E[\langle z_1,c\rangle\mid c]=As$ exactly, and $\sd[\langle z_1,c\rangle\mid c]=O(d^{-1/2})$.
\end{lemma}

\begin{proof}
The mean is \eqref{eq:m1} contracted with $c$. For the spread,
$\langle z_1,c\rangle=s\langle z_1,m\rangle+\beta\langle z_1,\xi\rangle$ with $\langle m,c\rangle=s$
deterministic; both $\langle z_1,m\rangle$ and the tangential components of a vMF variate
concentrate at rate $d^{-1/2}$.
\end{proof}

\begin{proposition}\label{prop:polarland}
$\law(\langle z_0,c\rangle\mid c)=\delta_{\cos j}$. By Corollary~\ref{cor:azimuth} the polar
dynamics $\dot\varphi=f(\varphi,t)$ do not depend on $\xi$, so all trajectories from
$\varphi_0=j$ solve the same scalar equation and $\law(\langle x_\infty,c\rangle\mid c)$ is a point
mass. In the limit $\sd[\langle z_1,c\rangle\mid c]\to0$ that point is $As$.
\end{proposition}

\begin{proof}
The first statement is immediate from Definition~\ref{def:model}. For the second, by
Lemma~\ref{lem:equiv} the field is $H_c$-equivariant and by Corollary~\ref{cor:azimuth} it is
confined to $\Pi$; all $\xi$ on the unit sphere of $c^{\perp}$ are related by elements of $H_c$,
so the induced scalar dynamics on $\varphi$ are the same for every $\xi$. A deterministic scalar
ODE from a fixed initial value has a single terminal value. The identification of that value with
$As$ is the content of Proposition~\ref{prop:sigmean} in the limit $\sigma\downarrow0$, and is
exact only in that limit.
\end{proof}

\begin{theorem}\label{thm:land}
In the limit of Proposition~\ref{prop:polarland},
\[
x_\infty=(As)\,c+\sqrt{1-A^{2}s^{2}}\ \xi,
\qquad
\E\langle x_\infty,z_1\rangle=A\Big(As^{2}+\beta\sqrt{1-A^{2}s^{2}}\Big).
\]
\end{theorem}

\begin{proof}
By Corollary~\ref{cor:azimuth}, $x_\infty=a\,c+b\,\xi$ with $a^{2}+b^{2}=1$ and azimuth $\xi$. By
Proposition~\ref{prop:polarland}, $a=As$, hence $b=\sqrt{1-A^{2}s^{2}}$ with the positive root
since $\varphi_t$ does not cross $0$. By the tower rule and \eqref{eq:m1},
\[
\E\langle x_\infty,z_1\rangle=\big\langle x_\infty,A(sc+\beta\xi)\big\rangle
=A\big(As^{2}+\beta\sqrt{1-A^{2}s^{2}}\big). \qedhere
\]
\end{proof}

\begin{corollary}\label{cor:neither}
The maximiser of $\E\langle\cdot,z_1\rangle$ given $c$ has $b=0$ (Remark~\ref{rem:whichopt}),
whereas $x_\infty$ has $b=\sqrt{1-A^{2}s^{2}}$. A draw from $\law(z_1\mid c)$ has random azimuth
and dispersed polar angle, whereas $x_\infty$ has azimuth determined by $z_0$ and a deterministic
polar angle.
\end{corollary}

\begin{measurement}\label{meas:land}
At $d=256$, $A=0.37$, $\beta=0.8$, $j=0.15$: Theorem~\ref{thm:land} predicts $0.3379$; the
population field attains $0.3360$; a network trained on $\mathcal M_1$ attains $0.3184$. For the
population field, $\E\langle x_\infty,c\rangle=0.1996$ against $As=0.2220$, and
$\sd\langle x_\infty,c\rangle=7.4\times10^{-8}$ against
$\sd\langle z_1,c\rangle=5.6\times10^{-2}$.
\end{measurement}

\begin{measurement}\label{meas:ring}
Let $\Pi_{\mathrm{cos}}:=\E\langle x_\infty,x_\infty'\rangle$ for two outputs generated from
independent draws of $\xi$ at fixed $c$. Corollary~\ref{cor:azimuth} and
Proposition~\ref{prop:polarland} give $\Pi_{\mathrm{cos}}=\E\langle x_\infty,c\rangle^{2}$. For the
population field, $\Pi_{\mathrm{cos}}=0.0400$ against $\E\langle x_\infty,c\rangle^{2}=0.0398$. For
a network trained on $\mathcal M_1$, $\Pi_{\mathrm{cos}}=0.1850$ against $0.0594$.
\end{measurement}

\subsection{Non-transfer of the pencil criterion}

\begin{proposition}\label{prop:codiag}
Conditional on $c$, the second moments $\Sigma_{0\mid c}$, $\Sigma_{1\mid c}$ and $C_{\mid c}$ of
the model of Definition~\ref{def:rho} lie in $\spn\{cc^{\T},P\}$ for every $\rho$. Consequently
$C_{\mid c}=C_{\mid c}^{\T}$, and $S_{\mid c}\in\spn\{\Sigma_{0\mid c},\Sigma_{1\mid c}\}$ whenever
that pencil is nondegenerate.
\end{proposition}

\begin{proof}
By Lemma~\ref{lem:invlaw} the conditional law is $H_c$-invariant, so each of the three matrices
commutes with the action of $H_c$. A symmetric matrix commuting with $SO(d-1)$ acting on
$c^{\perp}$ is a linear combination of $cc^{\T}$ and $P$; the same argument applied to
$C_{\mid c}+C_{\mid c}^{\T}$ and $C_{\mid c}-C_{\mid c}^{\T}$ shows the skew part is $H_c$-invariant
and skew, hence $0$, since no nonzero skew matrix on $\R^{d}$ commutes with $SO(d-1)$ for
$d\geq3$. Elements of $\spn\{cc^{\T},P\}$ commute pairwise, and three elements of a
$2$-dimensional space are linearly dependent, giving the last statement.
\end{proof}

\begin{measurement}\label{meas:codiag}
At the operating point, with $a:=\langle c,\Sigma c\rangle$ and $b:=\mathrm{tr}(P\Sigma)/(d-1)$,
\[
\|\Sigma_{0\mid c}-(acc^{\T}+bP)\|_{\max}=4.3\times10^{-6},
\quad
\|\Sigma_{1\mid c}-\cdot\|_{\max}=2.5\times10^{-4},
\quad
\|C_{\mid c}-\cdot\|_{\max}=3.8\times10^{-5},
\]
\[
\|C_{\mid c}-C_{\mid c}^{\T}\|_{\max}=4.8\times10^{-5},
\qquad
\|\Sigma_{0\mid c}\Sigma_{1\mid c}-\Sigma_{1\mid c}\Sigma_{0\mid c}\|_{\max}=1.4\times10^{-7}.
\]
\end{measurement}

\begin{remark}\label{rem:nottransfer}
Proposition~\ref{prop:codiag} establishes the algebraic conditions (i) and (ii) for the
conditional second moments. It does not establish the hypotheses of Theorem~\ref{thm:pencil},
which are Assumption~\ref{ass:flat}: jointly Gaussian endpoints, the linear interpolant
$x_t=(1-t)x_0+tx_1$, and $V(t)\succ0$ on $[0,1]$. The model of Definition~\ref{def:model} satisfies
none of the three: $z_1$ is von Mises--Fisher, the interpolant is the geodesic of
Definition~\ref{def:flow}, and by Lemma~\ref{lem:sigrank}, $\Var(\langle z_0,c\rangle\mid c)=0$.
Theorem~\ref{thm:pencil} therefore does not apply, and its conclusion fails.
\end{remark}

\begin{measurement}\label{meas:nontransfer}
Under the sweep of Definition~\ref{def:rho} with both conditional marginals fixed
(Lemma~\ref{lem:rhomarg}, verified: $\max_\rho|\E\langle z_1,c\rangle-\text{const}|=5.8\times10^{-4}$),
the population field's landing point moves. Values of $\E\langle x_\infty,c\rangle$:
\begin{center}
\begin{tabular}{cccccc}
\toprule
$\sigma\backslash\rho$ & $0$ & $0.25$ & $0.5$ & $0.75$ & $1$\\
\midrule
$0.005$ & $0.2237$ & $0.2199$ & $0.2198$ & $0.2066$ & $0.2019$\\
$0.01$  & $0.2301$ & $0.2224$ & $0.2184$ & $0.2142$ & $0.1987$\\
$0.10$  & $0.2671$ & $0.2619$ & $0.2560$ & $0.2549$ & $0.2515$\\
\bottomrule
\end{tabular}
\end{center}
Each row is strictly decreasing in $\rho$; under exchangeability a strict ordering of five values
has probability $1/60$. The maximum spread over $\rho$ at fixed $\sigma$ is $3.4\times10^{-2}$,
against a refit dispersion of $1.1\times10^{-2}$ measured from three independent refits at fixed
$(\sigma,\rho)$.
\end{measurement}

\begin{measurement}\label{meas:gainsphere}
Proposition~\ref{prop:sigmean} predicts
$\sd\langle x_\infty,c\rangle=\sd\langle z_1,c\rangle$ for $\sigma>0$. Measured at $\rho=1$, the
ratio $\sd\langle x_\infty,c\rangle/\sd\langle z_1,c\rangle$ has median $2.03$ over
$\sigma\in\{0.005,\dots,0.4\}$, ranging from $0.66$ to $3.11$. The predicted constancy of
$\E\langle x_\infty,c\rangle$ in $\sigma$ also fails: the value ranges from $0.1996$ at $\sigma=0$
to $0.2924$ at $\sigma=0.2$, against $As=0.2220$.
\end{measurement}

\begin{remark}\label{rem:attribution}
Measurements~\ref{meas:nontransfer} and \ref{meas:gainsphere} do not separate the contribution of
curvature from that of vMF non-Gaussianity. A flat non-Gaussian control with the same conditional
variance profile is required for that attribution and is not reported here.
\end{remark}

\section{The threshold}\label{sec:bstar}

\begin{definition}\label{def:delta}
$\Delta(\beta):=\E\langle x_\infty,z_1\rangle-\E\langle z_0,z_1\rangle$, i.e.\
Theorem~\ref{thm:land} minus \eqref{eq:src}, at fixed $(A,j)$. Let $\beta^{\star}(A,j)$ be the
first zero of $\Delta$ in $(0,1)$.
\end{definition}

\begin{proposition}\label{prop:bstar}
$\Delta(\beta)=0$ is equivalent, after dividing by $A>0$ and substituting $s=\sqrt{1-\beta^{2}}$, to
\begin{equation}\label{eq:bstar}
A(1-\beta^{2})+\beta\sqrt{1-A^{2}(1-\beta^{2})}=\sqrt{1-\beta^{2}}\,\cos j+\beta\sin j.
\end{equation}
The dimension enters only through $A=A_d(\kappa)$, so $\beta^{\star}$ is a function of $(A,j)$ alone.
\end{proposition}

\begin{proposition}\label{prop:exist}
$\Delta(0)=A(A-\cos j)$ and $\Delta(1)=A(1-\sin j)>0$ for $j\in[0,\pi/2)$. Hence
\[
\beta^{\star}\ \text{exists}\iff A<\cos j.
\]
\end{proposition}

\begin{proof}
At $\beta=0$, $s=1$: Theorem~\ref{thm:land} gives $A(A\cdot1+0)=A^{2}$ and \eqref{eq:src} gives
$A(\cos j+0)=A\cos j$. At $\beta=1$, $s=0$: Theorem~\ref{thm:land} gives $A(0+\sqrt{1-0})=A$ and
\eqref{eq:src} gives $A\sin j$. Since $\Delta$ is continuous with $\Delta(1)>0$, a zero in $(0,1)$
exists if and only if $\Delta(0)<0$.
\end{proof}

\begin{center}
\begin{tabular}{lcccccc}
\toprule
$\beta^{\star}$ & $j=0.00$ & $0.05$ & $0.15$ & $0.30$ & $0.50$ & $0.80$\\
\midrule
$A=0.15$ & 0.666 & 0.686 & 0.723 & 0.776 & 0.838 & 0.915\\
$A=0.25$ & 0.635 & 0.656 & 0.696 & 0.753 & 0.821 & 0.905\\
$A=0.37$ & 0.592 & 0.615 & 0.659 & 0.721 & 0.797 & 0.892\\
$A=0.50$ & 0.538 & 0.562 & 0.611 & 0.680 & 0.765 & 0.874\\
$A=0.70$ & 0.429 & 0.457 & 0.512 & 0.594 & 0.698 & ---\\
$A=0.90$ & 0.255 & 0.286 & 0.352 & 0.453 & --- & ---\\
\bottomrule
\end{tabular}
\end{center}
Each ``---'' satisfies $A\geq\cos j$, as Proposition~\ref{prop:exist} requires.

\begin{measurement}\label{meas:bstar}
Over the $27$ pairs $(A,j)$ with $j>0$ and $A<\cos j$, the zero of $\Delta$ located by linear
interpolation of the measured $\Delta$ on $\beta\in\{0,0.2,0.4,0.6,0.7,0.8,0.9,0.95\}$ agrees with
\eqref{eq:bstar} to a maximum absolute error of $3.3\times10^{-2}$.
\end{measurement}

\begin{remark}\label{rem:jzero}
At $j=0$ one has $z_0=c$ and $\Pi_{z_0}c=c-\langle c,c\rangle c=0$, so by
Proposition~\ref{prop:2plane} the field vanishes at $z_0$ and the trajectory is stationary. Then
$\Delta\equiv0$ and $\beta^{\star}$ is not identifiable from the flow, although \eqref{eq:bstar}
returns a value. The column $j=0$ of the table is therefore formal.
\end{remark}

\section{Starting at $z_0$ versus conditioning on $z_0$}\label{sec:cond}

\begin{proposition}\label{prop:det}
For the flat interpolant and $t>0$ the map $(z_0,z_t)\mapsto z_1$ is well defined:
\[
z_t=(1-t)z_0+tz_1\quad\Longrightarrow\quad z_1=\frac{z_t-(1-t)z_0}{t}.
\]
\end{proposition}

\begin{corollary}\label{cor:loss}
$\Var(\dot z_t\mid z_t,t,c,z_0)=0$ and $\Var(\dot z_t\mid z_t,t,c)>0$. The irreducible losses of
$\mathcal M_2$ and $\mathcal M_1$ are therefore $0$ and positive respectively, and the two
objectives of Definition~\ref{def:flow} are not comparable in value.
\end{corollary}

\begin{proof}
$\dot z_t=z_1-z_0$ is a deterministic function of $(z_0,z_t,t)$ by Proposition~\ref{prop:det};
$z_1$ is not determined by $(z_t,t,c)$.
\end{proof}

\begin{proposition}\label{prop:amp}
Let the auxiliary state be supplied at time $\alpha t$ with $\alpha\in[0,1)$. Then
\[
z_1=\Big(\frac{z_t}{1-t}-\frac{z_{\alpha t}}{1-\alpha t}\Big)\cdot\frac{(1-t)(1-\alpha t)}{t(1-\alpha)},
\qquad
\Big\|\frac{\partial z_1}{\partial z_t}\Big\|=\frac{1-\alpha t}{t(1-\alpha)}.
\]
\end{proposition}

\begin{proof}
From $z_t=(1-t)z_0+tz_1$ and $z_{\alpha t}=(1-\alpha t)z_0+\alpha t\,z_1$, divide by $(1-t)$ and
$(1-\alpha t)$ respectively and subtract to eliminate $z_0$:
\[
\frac{z_t}{1-t}-\frac{z_{\alpha t}}{1-\alpha t}
=z_1\Big(\frac{t}{1-t}-\frac{\alpha t}{1-\alpha t}\Big)
=z_1\cdot\frac{t\big[(1-\alpha t)-\alpha(1-t)\big]}{(1-t)(1-\alpha t)}
=z_1\cdot\frac{t(1-\alpha)}{(1-t)(1-\alpha t)},
\]
using $(1-\alpha t)-\alpha(1-t)=1-\alpha$. Solving for $z_1$ and differentiating in $z_t$ gives the
coefficient $\tfrac{1}{1-t}\cdot\tfrac{(1-t)(1-\alpha t)}{t(1-\alpha)}$.
\end{proof}

\begin{corollary}\label{cor:amp}
$\alpha=0$ gives $t^{-1}$, unbounded as $t\downarrow0$. At $(t,\alpha)=(0.05,0.9)$ the factor is
$\tfrac{1-0.045}{0.05\cdot0.1}=191$. The map of Proposition~\ref{prop:amp} exists for every
$\alpha<1$; the parameter $\alpha$ controls its conditioning, not its existence.
\end{corollary}

\begin{remark}\label{rem:origin}
At $t=0$ every pair $(z_0,z_1)$ with the given $z_0$ has $z_t=z_0$, so
Proposition~\ref{prop:det} fails and $\E[\dot z_0\mid z_0,c]$ averages over $z_1$. The population
optimum of $\mathcal M_2$ satisfies $\dot y=(y-z_0)/t$ in exact arithmetic, whose general solution
$y(t)=z_0+kt$ satisfies $y(0)=z_0$ for every constant $k$; the initial value problem at $t=0$ is
not well posed.
\end{remark}

\begin{measurement}\label{meas:m2}
At the operating point, networks trained to the two objectives of Definition~\ref{def:flow} at
matched compute attain $\E\langle x_\infty,z_1\rangle=0.3184$ for $\mathcal M_1$ and $0.2167$ for
$\mathcal M_2$, with terminal training losses $1.06\times10^{-3}$ and $5.19\times10^{-4}$. Under
maximal-update parameterisation with the learning rate tuned at width $256$ and transferred, five
seeds per cell, and checkpoint selection on a distributional criterion:
\begin{center}
\begin{tabular}{lccc}
\toprule
width & $\mathcal M_1$ & $\mathcal M_2$ & difference\\
\midrule
$256$  & $0.3138$ & $0.2153\ (\pm0.0127)$ & $0.0985$\\
$2048$ & $0.2904$ & $0.1864\ (\pm0.0107)$ & $0.1040$\\
\bottomrule
\end{tabular}
\end{center}
Over the same range the terminal loss of $\mathcal M_2$ falls from $4.52\times10^{-4}$ to
$2.56\times10^{-4}$ and its distributional error falls from $5.01\times10^{-2}$ to
$3.68\times10^{-2}$.
\end{measurement}

\begin{remark}\label{rem:claimstrength}
Proposition~\ref{prop:det} shows the coupling is recoverable on-path, so $\mathcal M_2$ is not
information-deficient. Proposition~\ref{prop:amp} and Remark~\ref{rem:origin} identify the two
defects: the recovery map is ill-conditioned as $t\downarrow0$, and the initial value problem is
degenerate at $t=0$. Measurement~\ref{meas:m2} shows the difference between the two objectives
does not decrease with capacity, while both the training loss and the distributional error of
$\mathcal M_2$ decrease.
\end{remark}

\begin{measurement}\label{meas:remedy}
Two interventions on $\mathcal M_2$ at the operating point. Truncating the integration to
$[\epsilon,1]$: the terminal value increases monotonically with $\epsilon$ over
$\epsilon\in\{0,0.01,0.02,0.05,0.1,0.2,0.35\}$, with Spearman correlation $0.96$ against
$\int_{\epsilon}^{1}t^{-1}dt$. Perturbing the auxiliary input by an independent tangential
rotation of magnitude $\sigma_{\mathrm{aux}}$ during training and sampling: the terminal value
rises from $0.2195$ at $\sigma_{\mathrm{aux}}=0$ to $0.3122$ at $\sigma_{\mathrm{aux}}=1$, against
$0.3195$ for $\mathcal M_1$.
\end{measurement}

\end{document}
